% math4cs-hw7-paths-cycles-solution.tex

% !TEX program = xelatex
%%%%%%%%%%%%%%%%%%%%
% see http://mirrors.concertpass.com/tex-archive/macros/latex/contrib/tufte-latex/sample-handout.pdf
% for how to use tufte-handout
\documentclass[a4paper, justified]{tufte-handout}

\input{hw-preamble} % feel free to modify this file if you understand LaTeX well
%%%%%%%%%%%%%%%%%%%%
\title{计算机数学 (8-paths-cycles: 路径与圈)}
\me{魏恒峰}{hfwei@hnu.edu.cn}{}{}
\date{2026年5月15日}
%%%%%%%%%%%%%%%%%%%%
\begin{document}
\maketitle
%%%%%%%%%%%%%%%%%%%%
\noplagiarism % PLEASE DON'T DELETE THIS LINE!
%%%%%%%%%%%%%%%%%%%%
\begin{abstract}
\end{abstract}
%%%%%%%%%%%%%%%%%%%%
\beginrequired
%%%%%%%%%%%%%%%

%%%%%%%%%%%%%%%
\begin{problem}
  设 $G = (V, E)$ 是无向图 (不一定是简单无向图),
  其中 $|E| = m$。
  请证明\footnote{这也说明了, $G$ 中度数为奇数的顶点数目为偶数。},
  \[
    \sum_{v \in V} \deg(v) = 2m.
  \]
\end{problem}

\begin{proof}
  每条边贡献了 2 度。所以~\footnote{这也被称为 ``握手定理''。},
  \[
    \sum_{v \in V} \deg(v) = 2m.
  \]
\end{proof}
%%%%%%%%%%%%%%%

%%%%%%%%%%%%%%%
\begin{problem}
  \begin{enumerate}[(1)]
    \item 请证明: 每个 $u\-v$ 闭道路 (closed walk) 都包含一个 $u\-v$ 路径 (path)。
    \item 请证明: 每个长度为奇数的闭道路 (closed walk) 都包含一个长度为奇数的圈 (cycle)~\footnote{
    ``长度''就是所含边的条数。}。
  \end{enumerate}

  \vspace{1em}
  \noindent (提示: 可用数学归纳法。如果你使用数学归纳法, 请注意数学归纳法的书写规范。)
\end{problem}

\begin{proof}
  \begin{enumerate}[(1)]
    \item 设 $W$ 是一条 $u\-v$ 闭道路。若 $W$ 中除起点/终点外没有重复顶点, 则 $W$ 本身就是一条 $u\-v$ 路径。

      否则, 设 $x$ 是 $W$ 中首次重复出现的顶点 (除 $u=v$ 外)。
      将 $W$ 在 $x$ 处拆成两条 $x\-x$ 闭道路 $W_1, W_2$。
      其中较短的一条 (或任一条) 仍可用归纳思想缩短:
      不断删除 ``来回折返'' 的形如 $\dots, a, b, a, \dots$ 的片段, 直到没有重复顶点为止。
      最终得到一条 $u\-v$ 路径。

    \item 对闭道路的长度 $l$ 作强数学归纳。
      \begin{description}
        \item[基础步骤:] $l = 1$。长度为 1 的闭道路即是长度为 1 的圈。显然成立。
        \item[归纳假设:] 假设每个长度为奇数 $1 \le l < k$ ($k$ 为奇数)
          的闭道路都包含一个长度为奇数的圈。
        \item[归纳步骤:] 考虑长度为奇数 $k > 1$ 的闭道路 $W$。
          分两种情况讨论:
          \begin{itemize}
            \item 如果 $W$ 中不包含重复顶点 (除了起点与终点), 则 $W$ 即为长度为奇数的圈。
            \item 不妨设 $W$ 中包含重复顶点 (除了起点与终点) $v$.
              $W$ 可看作两条起点、终点均为 $v$ 的闭道路。
              其中必有一条的长度为奇数, 记该闭道路为 $U$, 长度为 $u$,
              有 $1 \le u < k$。根据归纳假设, $U$ 中包含一个长度为奇数的圈。
          \end{itemize}
      \end{description}
  \end{enumerate}
\end{proof}
%%%%%%%%%%%%%%%

%%%%%%%%%%%%%%%
\begin{problem}
  设 $G$ 是一个简单无向图 (undirected simple graph) 且满足
  \[
    \delta(G) \ge k,
  \]
  其中 $k \in \mathbb{N}^{+}$ 为常数。
  请证明:
  \begin{enumerate}[(1)]
    \item $G$ 包含长度 $\ge k$ 的路径;
    \item 如果 $k \ge 2$, 则 $G$ 包含长度 $\ge k + 1$ 的圈。
  \end{enumerate}
  提示: By Extermality。
\end{problem}

\begin{proof}
  \begin{enumerate}[(1)]
    \item 考虑 $G$ 中的一条极大路径 $P$。设 $u$ 是 $P$ 的一个端点。
      因为 $P$ 是极大路径, $u$ 的所有邻居顶点都在 $P$ 上。
      因为 $\deg(u) \ge \delta(G) \ge k$ 且 $G$ 是简单图,
      所以除了端点 $u$, $P$ 至少还包含 $k$ 个顶点。
      故, $P$ 的长度 $\ge k$。
    \item 假设 $k \ge 2$。设 $v$ 是 $u$ 的邻居顶点中距离 $u$ 最远
      (在顺着路径 $P$ 的意义下) 的顶点。则边 $\set{u, v}$ 以及
      路径 $P$ 中从 $u$ 到 $v$ 的一段子路径构成了长度 $\ge k + 1$ 的圈。
  \end{enumerate}
\end{proof}
%%%%%%%%%%%%%%%

%%%%%%%%%%%%%%%
\begin{problem}
  考虑下图, 记为 $G$。
  \fig{width = 0.40\textwidth}{figs/Euler-graph}
  \begin{enumerate}[(1)]
    \item $G$ 是否是欧拉图? 请说明理由。
    \item 如果是欧拉图, 请将其分解为若干圈的组合, 并给出一个欧拉回路。
    　(可以自行为顶点编号, 也可以使用图上边的编号描述回路。)
  \end{enumerate}
\end{problem}

\begin{proof}
  \begin{enumerate}[(1)]
    \item $G$ 是欧拉图。因为 $G$ 中每个顶点的度数都是偶数。
    \item 可将 $G$ 分解成如下三个圈的组合:
      \begin{itemize}
        \item $A-B-F-H-K$
        \item $C-G-J$
        \item $D-E-I$
      \end{itemize}
      一个欧拉回路: $A-B-C-D-E-F-G-H-I-J-K$。
  \end{enumerate}
\end{proof}
%%%%%%%%%%%%%%%

%%%%%%%%%%%%%%%
\begin{problem}
  请判断下图 $G$ 是否是 Hamilton 图 (Hamiltonian graph), 并说明理由。
  \fig{width = 0.30\textwidth}{figs/Hamiltonian}
\end{problem}

\begin{proof}
  $G$ 是 Hamilton 图。

  该图可看作广义 Petersen 图 $GP(8, 3)$:
  外层 8 个顶点构成圈 $u_0u_1\cdots u_7u_0$,
  内层 8 个顶点通过 ``隔 3 连边'' 构成圈 $v_0v_3v_6v_1v_4v_7v_2v_5v_0$,
  且每个 $u_i$ 与 $v_i$ 相连。

  一个 Hamilton 圈为
  \[
    u_0, v_0, v_3, u_3, u_4, v_4, v_7, u_7, u_6, v_6, v_1, u_1, u_2, v_2, v_5, u_5, u_0.
  \]
  它经过每个顶点恰好一次并回到起点, 故 $G$ 是 Hamilton 图。
\end{proof}
%%%%%%%%%%%%%%%

%%%%%%%%%%%%%%%
\begin{problem}
  在圆周上取 $n$ 个点。
  将每个点与圆周上两个方向上各自距离它最近的两个点相连,
  得到一个 $4$-正则图 $G_n$。

  \noindent 请证明: 当 $n \ge 5$ 时, $G_n$ 可以表示为两个 Hamilton 圈的并。
\end{problem}

\begin{proof}
  将圆周上的点记为 $v_0, v_1, \dots, v_{n-1}$ (下标模 $n$)。
  则 $G_n$ 的边集为
  \[
    E(G_n) = \set{\set{v_i, v_{i+1}}, \set{v_i, v_{i+2}} \mid 0 \le i \le n-1}.
  \]
  换言之, $G_n = \mathrm{Circ}(n, \set{1, 2})$。

  先令
  \[
    H_1 = (v_0, v_1, v_2, \dots, v_{n-1}, v_0),
  \]
  它使用全部 ``距离为 1'' 的边, 因而是一条 Hamilton 圈。

  \begin{enumerate}[(1)]
    \item 若 $n$ 为奇数, 则 $\gcd(n, 2) = 1$。
      再令
      \[
        H_2 = (v_0, v_2, v_4, \dots, v_{n-1}, v_0),
      \]
      其中每一步都沿 ``距离为 2'' 的边前进。
      由于 $\gcd(n, 2) = 1$, 该序列依次访问全部 $n$ 个顶点并回到 $v_0$,
      故 $H_2$ 也是 Hamilton 圈。
      又 $H_2$ 只使用 ``距离为 2'' 的边, 与 $H_1$ 边不交, 故
      \[
        E(G_n) = E(H_1) \sqcup E(H_2).
      \]

      例如 $n = 5$ 时,
      \[
        H_1 = (v_0, v_1, v_2, v_3, v_4, v_0),
        \qquad
        H_2 = (v_0, v_2, v_4, v_1, v_3, v_0).
      \]

    \item 若 $n$ 为偶数, 则 $\gcd(n, 2) = 2$,
      ``距离为 2'' 的边只能组成两个长为 $n/2$ 的圈, 不能再单独取一条 Hamilton 圈。
      此时需把两类边混合使用。令
      \[
        H_1 = (v_0, v_1, \dots, v_{n-3}, v_{n-1}, v_{n-2}, v_0),
      \]
      \[
        H_2 = (v_0, v_2, v_4, \dots, v_{n-2}, v_{n-3}, v_{n-5}, \dots, v_1, v_{n-1}, v_0).
      \]
      两条闭链的每一步都是 ``距离为 1'' 或 ``距离为 2'' 的边,
      且各经过 $n$ 个不同顶点, 故 $H_1, H_2$ 都是 Hamilton 圈。
      例如 $n = 6$ 时,
      \[
        H_1 = (v_0, v_1, v_2, v_3, v_5, v_4, v_0),
        \qquad
        H_2 = (v_0, v_2, v_4, v_3, v_1, v_5, v_0).
      \]

      下面说明 $E(G_n) = E(H_1) \sqcup E(H_2)$。
      记
      \[
        E_1 = \set{\set{v_i, v_{i+1}} \mid 0 \le i \le n-1},
        \qquad
        E_2 = \set{\set{v_i, v_{i+2}} \mid 0 \le i \le n-1}.
      \]
      则 $E(G_n) = E_1 \sqcup E_2$, 且 $|E_1| = |E_2| = n$。

      \emph{第一步: $H_1$ 与 $H_2$ 边不交。}
      直接比较两条闭链可知, 它们没有公共边。

      \emph{第二步: $E(H_1) \cup E(H_2) = E(G_n)$。}
      在 $H_1$ 中, 除 $\set{v_{n-3}, v_{n-2}}$ 外,
      $E_1$ 中其余 $n-1$ 条边都出现一次;
      另外还出现 $E_2$ 中的
      $\set{v_{n-3}, v_{n-1}}$ 与 $\set{v_{n-2}, v_0}$。
      在 $H_2$ 中, 恰好补上 $\set{v_{n-3}, v_{n-2}}$ 与 $\set{v_{n-1}, v_0}$,
      并覆盖 $E_2$ 中其余 $n-2$ 条边。
      因此 $E(G_n) = E(H_1) \sqcup E(H_2)$。
  \end{enumerate}

  综上, 当 $n \ge 5$ 时, $G_n$ 都可表示为两个 Hamilton 圈的并。
\end{proof}
%%%%%%%%%%%%%%%

%%%%%%%%%%%%%%%%%%%%
% 如果没有需要订正的题目，可以把这部分删掉

\begincorrection
%%%%%%%%%%%%%%%%%%%%

%%%%%%%%%%%%%%%%%%%%
% 如果没有反馈，可以把这部分删掉
\beginfb

你可以写 (也可以发邮件)
\begin{itemize}
  \item 对课程及教师的建议与意见
  \item 教材中不理解的内容
  \item 希望深入了解的内容
  \item $\cdots$
\end{itemize}
%%%%%%%%%%%%%%%%%%%%
\end{document}
